% Fragmento público generado desde propietarios íntegros.
% Residencia: 52.9.
% Fuente: ../incoming/radion_monografia_20260827/manuscrito/sections/10_minkowski_hodge_lorentz.tex:1-55
\noindent L'holonomie orientée admet des réalisations causales et géométriques ultérieures dans Minkowski, Hodge, Lorentz et Cartan--Holst.\par\medskip

\paragraph{Construction de la métrique lorentzienne à partir du bloc APP}

\paragraph{Matrice APP de paramètres \(7\) et \(2\)}

La matrice centrale
\begin{equation}
B_c=\begin{pmatrix}7&2\\2&7\end{pmatrix}
=9P_++5P_-
\label{eq:Bc}
\end{equation}
est normalisée par son déterminant :
\begin{equation}
M_c=\frac{B_c}{\sqrt{45}}
=\exp\left(\frac12\log\frac95\,R\right)
\in SO^+(1,1).
\label{eq:Mc}
\end{equation}
La rapidité interne est
\[
\eta_c=\frac12\log\frac95.
\]
La publication projective associée satisfait
\begin{equation}
P_B^n=P_++\left(\frac59\right)^nP_-.
\label{eq:PB-contract}
\end{equation}
Le même rapport \(5/9\) qui apparaît comme contraction spectrale est reconnu ici
comme séparation d'un mode stable et d'un autre contractant. L'identité est
opératoire ; elle ne transforme pas toutes les occurrences de \(5/9\) en la même application.

\paragraph{Incidence \(6\to8\) et quotient de signature causale}

Soit \(M\) l'incidence typée entre six faces et huit sommets. La relation
spectrale récupérée est
\begin{equation}
MM^{\mathsf T}=12P_u+4P_-,
\label{eq:MMt}
\end{equation}
où \(P_u\) projette sur le mode uniforme et \(P_-\) sur le sous-espace
orienté. Le quotient survivant de dimension quatre se décompose comme
\begin{equation}
Q_4\simeq\R u\oplus E_-.
\label{eq:radion-Q4}
\end{equation}
En déclarant négatif le mode uniforme et positifs les trois modes orientés,
on obtient
\begin{equation}
B_{\mathrm{caus}}=\operatorname{diag}(-1,1,1,1).
\label{eq:minkowski}
\end{equation}
L'écran de Minkowski ne sélectionne pas l'état HMT. C'est la réalisation
causale ultérieure du quotient que l'incidence a déjà construit.

